\section{Stand der Technik}
\label{sec:StandDerTechnik}

Gegenwärtige Forschungsarbeiten und Publikationen evaluieren Betriebssysteme für eingebettete Systeme und das Internet of Things (IoT) vorrangig anhand architektonischer Eigenschaften sowie ihrer Ausführungsleistung. 

Eine umfassende qualitative Analyse verschiedener IoT-Betriebssysteme liefern \textit{Kanwal et al. (2023)}\footnote{\url{https://www.mdpi.com/2079-9292/12/3/708}}. Sie vergleichen unterschiedliche Betriebssysteme hinsichtlich Kernkriterien wie OS-Architektur, Programmiermodellen, Scheduling, Speicherverwaltung, Dateisystemen, sowie spezifischen Vor- und Nachteilen (Advantages und Limitations). Dabei bleibt die Betrachtung jedoch primär auf einer summarischen und konzeptionellen Ebene, ohne eine tiefergehende quantitative Bewertung der Implementierungs- und Konfigurationskomplexität vorzunehmen.

Darüber hinaus existieren Studien wie jene von \textit{Silva et al. (2019)}\footnote{\url{https://ieeexplore.ieee.org/document/8835860}}, die Contiki-NG, RIOT und Zephyr explizit auf ressourcenbeschränkten Geräten (Low-End Devices) vergleichen. Während hier ausführliche operative Benchmarks durchgeführt werden, liegt der Fokus in der Auswertung hauptsächlich auf der Verarbeitungszeit innerhalb des Netzwerkverkehrs und dem Verhalten der Betriebssysteme unter verschiedenen Lastbedingungen (Network Stack Benchmarks).

Zusätzlich finden sich in der technischen Praxis und industrienahen Publikationen (Grey Literature) Berichte, die spezifische Low-Level-Laufzeiteigenschaften vergleichen. Ein repräsentatives Beispiel hierfür ist die Publikation \textit{„Measuring Real-Time Operating System Performance Part II: Comparing FreeRTOS vs. Zephyr“}\footnote{\url{https://interrupt.memfault.com/blog/rtos-performance-freertos-vs-zephyr-part-2}}. Solche Analysen widmen sich den zeitlichen Verhaltensweisen von Echtzeitbetriebssystemen und vergleichen Metriken wie benötigte Taktzyklen, Interrupt-Latenzen oder Context-Switch-Zeiten. Eine ganzheitliche Gegenüberstellung, die den Entwicklungsaufwand, die Toolchain-Komplexität (beispielsweise durch Konzepte wie \texttt{Kconfig} und \texttt{DeviceTree} im Zephyr-Ökosystem) oder den Konfigurations-Overhead berücksichtigt, findet nicht statt.

\subsection{Zusammenfassung und Abgrenzung}
Zusammenfassend lässt sich feststellen, dass der aktuelle Stand der Technik Betriebssysteme zumeist detailliert miteinander vergleicht – sei es auf konzeptioneller Ebene oder durch Leistungsmessungen direkt auf dem Edge-Device. Eine pragmatische Analyse, die bewertet, wie hoch die initiale \textit{Komplexität} bei der konkreten Firmware-Erstellung ist, fehlt jedoch. Oftmals werden lediglich die reinen Taktzyklen und Latenzen einzelner Funktionen betrachtet, aber eine Analyse zur Verhältnismäßigkeit zwischen dem notwendigen Konfigurationsaufwand und der tatsächlichen Nutzlast (Features) steht noch aus. Genau hier setzt die vorliegende Arbeit an.
